The prime numbers look random: no formula produces them, twin primes keep appearing, and so do long gaps. Yet in bulk
they obey precise laws. Their number up to $x$ is close to $x/\log x$; they are shared equally among the residue classes
that can contain them; the number of prime factors of a typical integer follows the bell curve of chance. This book proves
these laws, with all the earlier volumes working together. It builds Dirichlet's $L$-functions to find primes in progressions, proves the prime number theorem, studies
the zeros of the zeta function and Riemann's explicit formula, in which the primes are a superposition of waves, one
for each zero, and explains the lead of the primes $4n+3$. The count of that lead up to $26861$ is proved here. The figure $99.59\%$ is quoted from Rubinstein and Sarnak, under the Riemann hypothesis and linear independence. Then come sieves and
twin primes, the Erd\H{o}s--Kac theorem, random models of the primes, exponential sums, and a first taste of modular
forms. Every numbered theorem is proved in full. A remark that quotes a deeper result, such as Deligne's bound or Maier's theorem, says so. Each chapter opens with the story behind its ideas.
